\section{\texorpdfstring{\sys}{WitnessRAG}}
\label{sec:method}

\sys treats memory retrieval as the search for facts that jointly satisfy
a question's information requirements. It operates in two stages
(\cref{fig:overview}). Memory construction organizes the interaction
history into a graph of dated propositions linked to their source
passages. Memory retrieval uses an evidence-grounded plan to search this
graph, verify candidate support, and assemble a compact context for the
reader. The planner specifies what to find; the executor establishes how
the stored facts connect.

\begin{figure*}[t]
  \centering
  \includegraphics[width=\textwidth]{figures/overview-standalone.pdf}
  \caption{\textbf{WitnessRAG architecture.}
  (1) \emph{Memory construction}: an extractor turns dialogue histories into propositions with triples,
  event dates, and source links. (2) \emph{Memory retrieval}: hybrid retrieval supplies evidence to
  a planner, which specifies conjunctive requirements, answer types, a
  temporal reference, and scoring weights. The executor joins graph facts
  under consistent bindings to produce witnesses. A conditional verifier
  checks candidates against source utterances; gaps and rejections feed
  bounded replanning. Within retrieval, accepted witnesses and candidate facts receive
  selection bonuses, while relevance ranking supplies additional evidence.
  The full profile delivers dated propositions and cached passage summaries
  to the reader. Blue outlines denote LLM operations, green outlines denote
  retrieval or deterministic processing, dotted links denote provenance,
  and dashed arrows denote feedback. Witness acceptance does not guarantee
  that all supporting facts survive the final context budget.}
  \Description{The architecture has two stages: memory construction and memory retrieval. Memory construction
  extracts propositions and indexes their time and importance metadata.
  Source passages and the dated proposition graph are connected by
  provenance. A question triggers evidence retrieval and an LLM planner.
  The planner specifies conjunctive requirements, answer type, temporal
  reference and weights. A graph executor finds consistent bindings without
  LLM calls. A conditional verifier checks source utterances, returning
  gaps or rejections to the planner. Accepted support and relevance-ranked
  candidates guide context selection. The reader receives dated statements
  and cached summaries and generates an answer.}
  \label{fig:overview}
\end{figure*}


\subsection{Memory Construction}
\label{sec:method-register}

The memory combines a passage store $\passages$ with a fact graph
$\mathcal G=(\mathcal N,\facts)$. Passages preserve the original dialogue
and its session structure. An \ac{LLM} extracts self-contained
propositions from the history, each represented by a statement and a
subject--relation--object triple. The statement preserves information
useful to the reader, while the triple provides the relational structure
used during search. Entity normalization connects recurring mentions
across passages, allowing facts recorded in different sessions to
participate in the same derivation.

Each fact retains a link to its source utterance and passage, an event
interval, and an importance score. Temporal expressions are interpreted
relative to the source session, distinguishing the time of an event
from the time at which it was mentioned. When no event time can be
resolved, the session date provides a reference. Importance approximates
the salience expressed in the source utterance through a
question-independent lexical signal.

The resulting memory is $\mem=(\passages,\facts)$. For a set of facts
$W$, let $\mathrm{src}(W)$ denote its source passages. These provenance
links connect graph-level reasoning to the original dialogue: a witness
can be traced back to the utterances from which its facts were extracted.
Memory construction is shared across all questions about the history.

\subsection{Evidence-Grounded Planning}
\label{sec:method-plan}

Retrieval begins with a hybrid search that combines dense and lexical
ranking through \ac{RRF}. The highest-ranked passages identify an initial
pool of evidence. An \ac{LLM} planner receives the question, relevant
dated facts from this pool, and the vocabulary of the memory. Planning
after retrieval lets the question be expressed using relations and
entities supported by the stored history.

The plan specifies a conjunctive query of the form introduced in
\cref{eq:cq}, together with the requested answer type, output form,
temporal reference, and ranking weights. Shared variables express
dependencies between facts: the value found by one atom constrains the
facts that can satisfy another. The answer may be an entity, a date, or
a collection of values. Type information helps rank candidate answers
and guides their subsequent verification. For comparison questions,
the graph retrieves the relevant premises and the reader performs the
comparison.

The temporal reference identifies the period against which memories
should be ranked. It may be the present, the beginning of the history,
or an interval named in the question. Together with the scoring weights,
this reference allows the same memory to be searched differently for
questions about recent events, first occurrences, and particular periods.
Once a plan is available, retrieval expands the evidence pool under
the corresponding ranking.

\subsection{Reference-Relative Ranking}
\label{sec:method-score}

Passage retrieval and witness search use a common scoring principle
that combines relevance, importance, and temporal proximity:
\begin{equation}
 s(x\mid q,\pi)=
 \lambda_s\,\mathrm{sim}(x)
 +\lambda_i\,\kappa_x
 +\lambda_t\,g(x;\rho).
 \label{eq:score}
\end{equation}
Here $\pi$ is the plan, $\kappa_x$ is the stored importance of item $x$,
and $\rho$ is the plan's temporal reference. The nonnegative weights sum
to one, with at least half of the weight assigned to similarity.
For a passage, $\mathrm{sim}(x)$ is its normalized hybrid retrieval
score. For a fact, it measures the match to a query atom using the
relation, entity arguments, and proposition text.

Temporal proximity decays with the distance between the item's event
interval $I_x$ and the reference interval:
\begin{equation}
 g(x;\rho)=\exp\!\left(-\frac{\delta(I_x,\rho)}{h_\rho}\right),
 \label{eq:proximity}
\end{equation}
where $\delta$ is zero for overlapping intervals and otherwise measures
their separation, and $h_\rho$ controls the decay scale. With the present
as reference, this term expresses recency; with an earlier interval, it
favors events close to that period. Time therefore acts as a
question-dependent preference rather than a universal preference for
recent memories. The same principle applies to passages and facts,
while their similarity terms reflect their different roles in retrieval.

\subsection{Witness Search}
\label{sec:method-prove}

The executor searches for witnesses of the planned query without
\ac{LLM} calls. Following \cref{sec:preliminaries}, a witness consists
of facts that instantiate every query atom under a consistent variable
assignment. Relation and entity matching accommodate differences
between the planner's language and the stored propositions. Once a
variable is bound, subsequent atoms must use the same entity, preventing
independently plausible facts from being combined into an inconsistent
answer.

Search proceeds through partial assignments using a bounded beam.
Atoms with more constants are considered first, and each accepted fact
extends the current assignment. For matched facts $f_1,\ldots,f_m$
forming a witness $W$, the search cost is
\begin{equation}
 J(W)=-\sum_{j=1}^{m}\log s(f_j\mid q,\pi)
       +\gamma\,|\mathrm{src}(W)|.
 \label{eq:witness-search}
\end{equation}
Each fact is scored against its corresponding atom. The first term
favors derivations with strong matches throughout the conjunction;
the second discourages unnecessary dispersion across passages.
The beam retains the lowest-cost partial assignments until all atoms
have been processed or no consistent extension remains.

For a single-answer plan, search first considers facts associated with
the retrieved evidence and expands to the full memory if necessary.
Set-valued plans search across the memory and retain support for
individual members. Candidate acceptance also considers raw matching
strength, extraction confidence, and the selectivity of the result,
so temporal or importance bonuses alone cannot make a weak match
acceptable. Single-answer plans require a selective result; set-valued
plans assess support separately for each candidate.

When search fails, it identifies an unsatisfied atom together with the
bindings established so far. This gap describes which relation remains
unresolved and provides a concrete target for revising the plan.

\subsection{Verification and Replanning}
\label{sec:method-verify}

A conditional verifier checks candidate answers against their source
utterances, using the question, plan, and dated supporting facts.
It assesses whether the evidence concerns the intended entity, period,
and answer type, and whether the source dialogue supports the proposed
derivation. Set members are assessed individually, allowing a subset
of candidates to be retained.

Verification is invoked when required by the evidence-delivery policy;
support already covered by the retrieved passages can be accepted
without an additional check. If no candidate is accepted, execution
gaps and verification feedback guide a new plan. The planner can revise
a relation, an entity, or the temporal interpretation, after which the
executor searches again. The cycle ends on acceptance or when its
search budget is exhausted. Relevance-based retrieval supplies evidence
when no witness is accepted.

\begin{algorithm}[t]
\caption{Witness-guided memory retrieval}
\label{alg:search}
\small
\begin{algorithmic}[1]
\Require Question $q$, memory $\mem$, cycle limit $R$, context budget $b$
\State $E\gets\search(q,\mem)$; $\Delta\gets\emptyset$
\State $W^+\gets\emptyset$
\For{$c=1,\ldots,R$}
  \State $\pi\gets\plan(q,E,\Delta)$
  \State $E\gets E\cup\search(q,\mem;\pi)$
  \State $(W,\Delta)\gets\prove(\pi,\mem,E)$
  \State $(W^+,\Delta)\gets\textsc{Assess}(q,\pi,W,\Delta)$
    \Comment{conditional verification}
  \If{$W^+\neq\emptyset$} \State \textbf{break} \EndIf
\EndFor
\State $\ctx\gets\textsc{Assemble}(q,\mem,\pi,W^+,b)$
\State \Return $\mathrm{Read}(q,\ctx)$
\end{algorithmic}
\end{algorithm}

\subsection{Context Assembly and Cost}
\label{sec:method-respond}
\label{sec:method-guarantees}

The final context combines dated propositions with summaries of their
source passages. Selection combines a preference for facts belonging to accepted witnesses
with candidate support and relevance to the question and the plan's atoms. Deduplication and limits on contributions from a
single utterance encourage coverage within the context budget.
The selected propositions are organized by session and annotated with
their event dates. Passage summaries supply surrounding information
that individual propositions may omit; they are generated independently
of the question and reused across retrievals.

Witness support thus guides context construction without excluding
additional relevant evidence. The reader uses this context to formulate
the answer and perform any required comparison or date arithmetic.
The distinction between finding support and compressing it into the
reader's context is discussed in \cref{sec:discussion}.

For at most $R$ retrieval cycles, the controller makes at most $R$
planning calls and $R$ verification calls, followed by one reader call.
Graph execution requires no \ac{LLM} calls. Extraction and passage
summarization contribute additional costs that can be amortized across
questions on the same history. This organization concentrates language
model use on interpretation, source assessment, and answer generation,
while relational composition is handled by the executor.
